% Level 02 - Inference service components
% Question answered: how does a scene become candidate plume objects?
% Facts: report D3 (L1C, B11/B12, target + up to 6 references), NFR1 (15 minutes per 10 x 10 km granule).
% Derived: a chip is 100 x 100 pixels at 20 m = 2 x 2 km, so a 10 x 10 km granule is 25 chips.
\documentclass[tikz,border=10pt]{standalone}
% =============================================================================
% MethaneWatch diagram style  (shared by every diagram under diagrams/)
% -----------------------------------------------------------------------------
% Why one shared file: every diagram must look like part of the same set, and a
% colour change (for example a layer re-coloured) should be one edit, not 33.
%
% Usage, at the top of any diagram:
%   \documentclass[tikz,border=10pt]{standalone}
%   % =============================================================================
% MethaneWatch diagram style  (shared by every diagram under diagrams/)
% -----------------------------------------------------------------------------
% Why one shared file: every diagram must look like part of the same set, and a
% colour change (for example a layer re-coloured) should be one edit, not 33.
%
% Usage, at the top of any diagram:
%   \documentclass[tikz,border=10pt]{standalone}
%   % =============================================================================
% MethaneWatch diagram style  (shared by every diagram under diagrams/)
% -----------------------------------------------------------------------------
% Why one shared file: every diagram must look like part of the same set, and a
% colour change (for example a layer re-coloured) should be one edit, not 33.
%
% Usage, at the top of any diagram:
%   \documentclass[tikz,border=10pt]{standalone}
%   \input{../_common/mwstyle.tex}
%   \begin{document} ... \end{document}
%
% Colour meaning is fixed across the whole set (see diagrams/README.md). Two colours plus grey:
%   blue = MethaneWatch's own parts (application, data, models, science, services)
%   red  = evidence, security, audit and hard limits (the things that must never be bypassed)
%   grey = people, cloud and infrastructure, external systems, notes
% The older names green, orange, purple, teal and gold are kept as aliases so that no diagram
% source needs to change; they now resolve to blue (green, orange, purple) or grey (teal, gold).
% =============================================================================
\usepackage[T1]{fontenc}
\usepackage[scaled=0.95]{helvet}
\renewcommand{\familydefault}{\sfdefault}
\hyphenpenalty=10000 \exhyphenpenalty=10000 % never hyphenate words inside diagrams
\usepackage{pifont}
\usepackage{amssymb}
\usepackage{array}
\usepackage{booktabs}
\usetikzlibrary{positioning,arrows.meta,shapes.geometric,shapes.multipart,%
  shapes.misc,fit,backgrounds,calc,decorations.pathreplacing,matrix,patterns}

% ---- palette ----------------------------------------------------------------
\definecolor{mwink}{HTML}{1F2933}
\definecolor{mwgrey}{HTML}{6B7785}
\definecolor{mwline}{HTML}{4A5560}
\definecolor{mwlight}{HTML}{F2F4F7}
\definecolor{mwblue}{HTML}{1F6FB2}   \definecolor{mwblueL}{HTML}{E3EFFA}
\definecolor{mwred}{HTML}{B5483F}    \definecolor{mwredL}{HTML}{F8E3E0}
% aliases (see note above)
\colorlet{mwgreen}{mwblue}   \colorlet{mwgreenL}{mwblueL}
\colorlet{mworange}{mwblue}  \colorlet{mworangeL}{mwblueL}
\colorlet{mwpurple}{mwblue}  \colorlet{mwpurpleL}{mwblueL}
\colorlet{mwteal}{mwgrey}    \colorlet{mwtealL}{mwlight}
\colorlet{mwgold}{mwgrey}    \colorlet{mwgoldL}{mwlight}

% ---- node and edge styles ---------------------------------------------------
\tikzset{
  >={Stealth[length=2.2mm]},
  every node/.style={font=\sffamily\small,text=mwink},
  box/.style={rectangle,rounded corners=2.5pt,draw=mwline,line width=0.7pt,
    fill=white,align=center,inner sep=5pt,minimum height=7mm},
  blue/.style={draw=mwblue,fill=mwblueL},
  green/.style={draw=mwgreen,fill=mwgreenL},
  orange/.style={draw=mworange,fill=mworangeL},
  purple/.style={draw=mwpurple,fill=mwpurpleL},
  red/.style={draw=mwred,fill=mwredL},
  teal/.style={draw=mwteal,fill=mwtealL},
  gold/.style={draw=mwgold,fill=mwgoldL},
  grey/.style={draw=mwgrey,fill=mwlight},
  strong/.style={line width=1.6pt},
  ext/.style={box,dashed,draw=mwgrey,fill=mwlight},
  actor/.style={box,gold,rounded corners=8pt,minimum height=9mm},
  store/.style={cylinder,shape border rotate=90,aspect=0.22,draw=mwgreen,fill=mwgreenL,
    line width=0.7pt,align=center,inner sep=3pt,minimum height=11mm,minimum width=20mm},
  group/.style={rectangle,rounded corners=4pt,draw=mwline,dashed,line width=0.7pt,
    inner sep=7pt},
  glabel/.style={anchor=north west,font=\sffamily\bfseries\footnotesize,text=mwgrey,inner sep=2pt},
  note/.style={font=\sffamily\footnotesize,text=mwgrey,align=left,inner sep=1pt},
  tiny/.style={font=\sffamily\scriptsize},
  head/.style={font=\sffamily\bfseries\small},
  flow/.style={->,line width=0.8pt,draw=mwline},
  dflow/.style={->,line width=0.8pt,draw=mwline,dashed},
  bflow/.style={<->,line width=0.8pt,draw=mwline},
  lab/.style={font=\sffamily\scriptsize,text=mwgrey,fill=white,inner sep=1pt},
}

% ---- title / footer helpers -------------------------------------------------
% \dtitle{Finding or name}{One-line subtitle: what the diagram answers}
\newcommand{\dtitle}[2]{%
  \node[anchor=south west,font=\sffamily\bfseries\Large,text=mwink,inner sep=0pt] (dt)
    at ([yshift=12mm]current bounding box.north west) {#1};
  \node[anchor=north west,font=\sffamily\footnotesize,text=mwgrey,inner sep=0pt]
    at ([yshift=-1.2mm]dt.south west) {#2};}
% \dfoot{Source or reading note}
\newcommand{\dfoot}[2][14cm]{%
  \node[anchor=north west,note,text width=#1] at
    ([yshift=-4mm]current bounding box.south west) {#2};}
% in-node text helpers: \nm{Name} bold name, \tech{technology} small grey italic
\newcommand{\nm}[1]{{\bfseries\small #1}}
\newcommand{\tech}[1]{{\scriptsize\itshape\color{mwgrey} #1}}
% check mark used in matrices
\newcommand{\yes}{\ding{51}}

%   \begin{document} ... \end{document}
%
% Colour meaning is fixed across the whole set (see diagrams/README.md). Two colours plus grey:
%   blue = MethaneWatch's own parts (application, data, models, science, services)
%   red  = evidence, security, audit and hard limits (the things that must never be bypassed)
%   grey = people, cloud and infrastructure, external systems, notes
% The older names green, orange, purple, teal and gold are kept as aliases so that no diagram
% source needs to change; they now resolve to blue (green, orange, purple) or grey (teal, gold).
% =============================================================================
\usepackage[T1]{fontenc}
\usepackage[scaled=0.95]{helvet}
\renewcommand{\familydefault}{\sfdefault}
\hyphenpenalty=10000 \exhyphenpenalty=10000 % never hyphenate words inside diagrams
\usepackage{pifont}
\usepackage{amssymb}
\usepackage{array}
\usepackage{booktabs}
\usetikzlibrary{positioning,arrows.meta,shapes.geometric,shapes.multipart,%
  shapes.misc,fit,backgrounds,calc,decorations.pathreplacing,matrix,patterns}

% ---- palette ----------------------------------------------------------------
\definecolor{mwink}{HTML}{1F2933}
\definecolor{mwgrey}{HTML}{6B7785}
\definecolor{mwline}{HTML}{4A5560}
\definecolor{mwlight}{HTML}{F2F4F7}
\definecolor{mwblue}{HTML}{1F6FB2}   \definecolor{mwblueL}{HTML}{E3EFFA}
\definecolor{mwred}{HTML}{B5483F}    \definecolor{mwredL}{HTML}{F8E3E0}
% aliases (see note above)
\colorlet{mwgreen}{mwblue}   \colorlet{mwgreenL}{mwblueL}
\colorlet{mworange}{mwblue}  \colorlet{mworangeL}{mwblueL}
\colorlet{mwpurple}{mwblue}  \colorlet{mwpurpleL}{mwblueL}
\colorlet{mwteal}{mwgrey}    \colorlet{mwtealL}{mwlight}
\colorlet{mwgold}{mwgrey}    \colorlet{mwgoldL}{mwlight}

% ---- node and edge styles ---------------------------------------------------
\tikzset{
  >={Stealth[length=2.2mm]},
  every node/.style={font=\sffamily\small,text=mwink},
  box/.style={rectangle,rounded corners=2.5pt,draw=mwline,line width=0.7pt,
    fill=white,align=center,inner sep=5pt,minimum height=7mm},
  blue/.style={draw=mwblue,fill=mwblueL},
  green/.style={draw=mwgreen,fill=mwgreenL},
  orange/.style={draw=mworange,fill=mworangeL},
  purple/.style={draw=mwpurple,fill=mwpurpleL},
  red/.style={draw=mwred,fill=mwredL},
  teal/.style={draw=mwteal,fill=mwtealL},
  gold/.style={draw=mwgold,fill=mwgoldL},
  grey/.style={draw=mwgrey,fill=mwlight},
  strong/.style={line width=1.6pt},
  ext/.style={box,dashed,draw=mwgrey,fill=mwlight},
  actor/.style={box,gold,rounded corners=8pt,minimum height=9mm},
  store/.style={cylinder,shape border rotate=90,aspect=0.22,draw=mwgreen,fill=mwgreenL,
    line width=0.7pt,align=center,inner sep=3pt,minimum height=11mm,minimum width=20mm},
  group/.style={rectangle,rounded corners=4pt,draw=mwline,dashed,line width=0.7pt,
    inner sep=7pt},
  glabel/.style={anchor=north west,font=\sffamily\bfseries\footnotesize,text=mwgrey,inner sep=2pt},
  note/.style={font=\sffamily\footnotesize,text=mwgrey,align=left,inner sep=1pt},
  tiny/.style={font=\sffamily\scriptsize},
  head/.style={font=\sffamily\bfseries\small},
  flow/.style={->,line width=0.8pt,draw=mwline},
  dflow/.style={->,line width=0.8pt,draw=mwline,dashed},
  bflow/.style={<->,line width=0.8pt,draw=mwline},
  lab/.style={font=\sffamily\scriptsize,text=mwgrey,fill=white,inner sep=1pt},
}

% ---- title / footer helpers -------------------------------------------------
% \dtitle{Finding or name}{One-line subtitle: what the diagram answers}
\newcommand{\dtitle}[2]{%
  \node[anchor=south west,font=\sffamily\bfseries\Large,text=mwink,inner sep=0pt] (dt)
    at ([yshift=12mm]current bounding box.north west) {#1};
  \node[anchor=north west,font=\sffamily\footnotesize,text=mwgrey,inner sep=0pt]
    at ([yshift=-1.2mm]dt.south west) {#2};}
% \dfoot{Source or reading note}
\newcommand{\dfoot}[2][14cm]{%
  \node[anchor=north west,note,text width=#1] at
    ([yshift=-4mm]current bounding box.south west) {#2};}
% in-node text helpers: \nm{Name} bold name, \tech{technology} small grey italic
\newcommand{\nm}[1]{{\bfseries\small #1}}
\newcommand{\tech}[1]{{\scriptsize\itshape\color{mwgrey} #1}}
% check mark used in matrices
\newcommand{\yes}{\ding{51}}

%   \begin{document} ... \end{document}
%
% Colour meaning is fixed across the whole set (see diagrams/README.md). Two colours plus grey:
%   blue = MethaneWatch's own parts (application, data, models, science, services)
%   red  = evidence, security, audit and hard limits (the things that must never be bypassed)
%   grey = people, cloud and infrastructure, external systems, notes
% The older names green, orange, purple, teal and gold are kept as aliases so that no diagram
% source needs to change; they now resolve to blue (green, orange, purple) or grey (teal, gold).
% =============================================================================
\usepackage[T1]{fontenc}
\usepackage[scaled=0.95]{helvet}
\renewcommand{\familydefault}{\sfdefault}
\hyphenpenalty=10000 \exhyphenpenalty=10000 % never hyphenate words inside diagrams
\usepackage{pifont}
\usepackage{amssymb}
\usepackage{array}
\usepackage{booktabs}
\usetikzlibrary{positioning,arrows.meta,shapes.geometric,shapes.multipart,%
  shapes.misc,fit,backgrounds,calc,decorations.pathreplacing,matrix,patterns}

% ---- palette ----------------------------------------------------------------
\definecolor{mwink}{HTML}{1F2933}
\definecolor{mwgrey}{HTML}{6B7785}
\definecolor{mwline}{HTML}{4A5560}
\definecolor{mwlight}{HTML}{F2F4F7}
\definecolor{mwblue}{HTML}{1F6FB2}   \definecolor{mwblueL}{HTML}{E3EFFA}
\definecolor{mwred}{HTML}{B5483F}    \definecolor{mwredL}{HTML}{F8E3E0}
% aliases (see note above)
\colorlet{mwgreen}{mwblue}   \colorlet{mwgreenL}{mwblueL}
\colorlet{mworange}{mwblue}  \colorlet{mworangeL}{mwblueL}
\colorlet{mwpurple}{mwblue}  \colorlet{mwpurpleL}{mwblueL}
\colorlet{mwteal}{mwgrey}    \colorlet{mwtealL}{mwlight}
\colorlet{mwgold}{mwgrey}    \colorlet{mwgoldL}{mwlight}

% ---- node and edge styles ---------------------------------------------------
\tikzset{
  >={Stealth[length=2.2mm]},
  every node/.style={font=\sffamily\small,text=mwink},
  box/.style={rectangle,rounded corners=2.5pt,draw=mwline,line width=0.7pt,
    fill=white,align=center,inner sep=5pt,minimum height=7mm},
  blue/.style={draw=mwblue,fill=mwblueL},
  green/.style={draw=mwgreen,fill=mwgreenL},
  orange/.style={draw=mworange,fill=mworangeL},
  purple/.style={draw=mwpurple,fill=mwpurpleL},
  red/.style={draw=mwred,fill=mwredL},
  teal/.style={draw=mwteal,fill=mwtealL},
  gold/.style={draw=mwgold,fill=mwgoldL},
  grey/.style={draw=mwgrey,fill=mwlight},
  strong/.style={line width=1.6pt},
  ext/.style={box,dashed,draw=mwgrey,fill=mwlight},
  actor/.style={box,gold,rounded corners=8pt,minimum height=9mm},
  store/.style={cylinder,shape border rotate=90,aspect=0.22,draw=mwgreen,fill=mwgreenL,
    line width=0.7pt,align=center,inner sep=3pt,minimum height=11mm,minimum width=20mm},
  group/.style={rectangle,rounded corners=4pt,draw=mwline,dashed,line width=0.7pt,
    inner sep=7pt},
  glabel/.style={anchor=north west,font=\sffamily\bfseries\footnotesize,text=mwgrey,inner sep=2pt},
  note/.style={font=\sffamily\footnotesize,text=mwgrey,align=left,inner sep=1pt},
  tiny/.style={font=\sffamily\scriptsize},
  head/.style={font=\sffamily\bfseries\small},
  flow/.style={->,line width=0.8pt,draw=mwline},
  dflow/.style={->,line width=0.8pt,draw=mwline,dashed},
  bflow/.style={<->,line width=0.8pt,draw=mwline},
  lab/.style={font=\sffamily\scriptsize,text=mwgrey,fill=white,inner sep=1pt},
}

% ---- title / footer helpers -------------------------------------------------
% \dtitle{Finding or name}{One-line subtitle: what the diagram answers}
\newcommand{\dtitle}[2]{%
  \node[anchor=south west,font=\sffamily\bfseries\Large,text=mwink,inner sep=0pt] (dt)
    at ([yshift=12mm]current bounding box.north west) {#1};
  \node[anchor=north west,font=\sffamily\footnotesize,text=mwgrey,inner sep=0pt]
    at ([yshift=-1.2mm]dt.south west) {#2};}
% \dfoot{Source or reading note}
\newcommand{\dfoot}[2][14cm]{%
  \node[anchor=north west,note,text width=#1] at
    ([yshift=-4mm]current bounding box.south west) {#2};}
% in-node text helpers: \nm{Name} bold name, \tech{technology} small grey italic
\newcommand{\nm}[1]{{\bfseries\small #1}}
\newcommand{\tech}[1]{{\scriptsize\itshape\color{mwgrey} #1}}
% check mark used in matrices
\newcommand{\yes}{\ding{51}}

\begin{document}
\begin{tikzpicture}[
  c/.style={box,text width=3.0cm,align=left,inner sep=4pt,minimum height=1.9cm,font=\sffamily\footnotesize}]

  \node[c,green] (in) at (0,0) {\nm{Scene stack}\\ One target scene and up to six reference scenes of the same site, B11 and B12 at 20 m};
  \node[c,green] (chip) at (4.0,0) {\nm{inference/tiling}\\ Cut the granule into 100 x 100 chips (25 for 10 x 10 km); build ref\_valid};
  \node[c,orange] (pre) at (8.0,0) {\nm{inference/preprocess}\\ Same normalisation as training; shape and band checks};
  \node[c,orange,strong] (pred) at (12.0,0) {\nm{inference/predictor}\\\tech{approved model}\\ Forward pass, per-pixel plume probability};
  \node[c,orange] (thr) at (12.0,-3.4) {\nm{inference/postprocess}\\ Threshold fixed on validation; drop objects below minimum size};
  \node[c,orange] (cc) at (8.0,-3.4) {\nm{connected components}\\ One object per plume candidate; merge across chip borders};
  \node[c,blue] (poly) at (4.0,-3.4) {\nm{inference/polygonise}\\ Georeferenced polygons (GeoJSON) and raster overlay};
  \node[c,blue,strong] (out) at (0,-3.4) {\nm{Candidate objects}\\ Mask URI, confidence, bounding box, model version, scene id};

  \node[c,grey,dashed,minimum height=1.5cm] (reg) at (12.0,3.1) {\nm{Model registry}\\ Loaded once per worker; version stored on every result};
  \node[c,purple,minimum height=1.5cm] (next) at (-4.0,-3.4) {\nm{Science chain}\\ verification, IME, attribution (see science-chain)};
  \node[c,red,minimum height=1.5cm,text width=3.4cm] (bad) at (4.0,3.1) {\nm{Unusable input}\\ Clouds or invalid pixels: no prediction is trusted; goes to review as Unreviewable};

  \draw[flow] (in) -- (chip);
  \draw[flow] (chip) -- (pre);
  \draw[flow] (pre) -- (pred);
  \draw[flow] (pred) -- (thr);
  \draw[flow] (thr) -- (cc);
  \draw[flow] (cc) -- (poly);
  \draw[flow] (poly) -- (out);
  \draw[flow] (out) -- (next);
  \draw[dflow] (reg) -- (pred);
  \draw[dflow] (chip.north) -- (bad.south);

  \dtitle{A granule is 25 chips; every chip carries its own validity mask}
         {Level 02 - inference service. Same code path locally and in the Beta GPU worker; only the model file and the adapters change.}
  \dfoot[17cm]{Target budget (NFR1): ingestion, preprocessing and inference for one 10 x 10 km granule within 15 minutes on a standard GPU instance, measured once all inputs are available.
    The 24 h figure is separate: satellite pass to candidate mask, a Beta target.}
\end{tikzpicture}
\end{document}
